% CHECKPOINT34: isolated new source; all checkpoint33 sources are unchanged.
\statement{A.107}{Proved}{The relative central tower and the actual modular quotient automorphism}
Retain A.103--A.106, and define the closed relative central series
\[
 N_0=N,\qquad N_{j+1}=\overline{[N_j,F]},\qquad
 N_\infty=\bigcap_{j\geq0}N_j .
 \tag{A.107.1}
\]
The actual normalized endomorphism
\[
 \psi_d(x)=x^d,\qquad \psi_d(y)=u^{-1}y^du
 \tag{A.107.2}
\]
induces automorphisms of every \(F/N_j\) and of \(F/N_\infty\).
On each of these actual quotients it conjugates the original
complex conjugation \(c\) to \(\alpha\). Moreover
\[
 \psi_d\text{ is onto }F
 \quad\Longleftrightarrow\quad
 \psi_d|_{N_\infty}:N_\infty\to N_\infty\text{ is onto}.
 \tag{A.107.3}
\]
For the original modular marking
\[
 B=C_2(s)\amalg C_3(r),\quad
 F=\ker(B\to S_3),\quad
 \delta=sr,\quad x=\delta^{-2},\quad y=sxs,
\]
the subgroup \(N_1\) is normal in \(B\). On the actual quotient
\(B_1=B/N_1\), the assignments
\[
 \Psi_1(s)=v=u^{-1}s=su,\qquad
 \Psi_1(\delta)=\delta^d
 \tag{A.107.4}
\]
define an automorphism whose restriction to \(F/N_1\) is (A.107.2)
and whose quotient map on \(S_3\) is the identity. They give
\[
 \alpha|_{B_1}=\Psi_1 c_B\Psi_1^{-1},
 \tag{A.107.5}
\]
where \(c_B(s)=s\), \(c_B(\delta)=\delta^{-1}\).
The assertion on \(B_1\) does not assert such an extension to \(B\)
or to every \(B/N_j\).

\begin{proof}
\emph{The actual geometric stability.}
Write \(\rho=\operatorname{Ad}(r)|_F\), so that \(\rho\) cycles
\(x,y,z=(xy)^{-1}\), and \(\iota=\operatorname{Ad}(s)|_F\), which swaps
\(x,y\). The signed integral matrices
\[
 J=\begin{pmatrix}0&-1\\1&-1\end{pmatrix},\qquad
 S_0=\begin{pmatrix}0&-1\\1&0\end{pmatrix}
\]
satisfy, for every \(g\in F\),
\[
 \begin{gathered}
 q(\rho(g))=Jq(g)J^{-1}(-I)^{\operatorname{exp}_y(g)},\\
 q(\iota(g))=S_0q(g)S_0^{-1}.
 \end{gathered}
 \tag{A.107.6}
\]
The exponent in the first line is read modulo two. These identities
follow on the free generators, hence on all of \(F\).
The actual modulo-four images
\(X_4=I+2E_{12}\), \(Y_4=I-2E_{21}\) are independent commuting
involutions. More explicitly,
\[
 q_4(g)=I+2\bigl(\operatorname{exp}_x(g)E_{12}
                       -\operatorname{exp}_y(g)E_{21}\bigr).
 \tag{A.107.7}
\]
Consequently every \(n\in N\) has both exponent parities zero.
Its sign in (A.107.6) is one; thus \(\rho(N)\subseteq N\), and
\(\rho^3=1\) gives equality. The second line gives
\(\iota(N)=N\). Since \(N\) is already normal in \(F\), it is normal
in \(B\). Closed-commutator functoriality proves the same assertion
for every \(N_j\) and \(N_\infty\).
This uses the full signed kernel. For comparison,
\(y^3\in\ker(q_3)\), while \(q_3(\rho(y^3))=-I\);
but \(y^3\notin N\) by (A.107.7). No characteristicity under all
automorphisms of \(F\) is asserted.

\emph{Automorphy on the whole relative central tower.}
A.104 proves that (A.107.2) induces an automorphism of \(C=F/N\),
that \(\psi_d^{-1}(N)=N\), and that both \(\alpha,c\) preserve \(N\).
Functoriality makes \(\psi_d\) preserve every \(N_j\), while
\(\alpha,c\) preserve each of them with equality.
Induct on \(j\geq1\). Assume that \(\psi_d\) is an automorphism on
\(F/N_{j-1}\), and put
\[
 E_j=F/N_j,\qquad Z_j=N_{j-1}/N_j,\qquad
 L_j=\psi_d(E_j).
\]
Here \(Z_j\) is central in \(E_j\), and the induction hypothesis gives
\(E_j=Z_jL_j\). Thus \(E_j'\subseteq L_j\): commutators of elements
of \(Z_jL_j\) are commutators of elements of \(L_j\), and \(L_j\)
is compact, hence closed. Since \(N_j\subseteq F'\),
\((E_j)_{\rm ab}=F_{\rm ab}\). The action of \(\psi_d\) on this
abelianization is multiplication by the unit \(d\), so
\(L_jE_j'=E_j\). It follows that \(L_j=E_j\).

The finitely generated profinite group \(E_j\) is Hopfian.
For completeness, it has only finitely many open subgroups of each
fixed index; the usual finite-index intersection argument therefore
makes a surjective endomorphism injective. Hence the induced map
on \(E_j\) is an automorphism. This completes the induction.
In particular \(\psi_d^{-1}(N_j)=N_j\) for all \(j\).

Compactness identifies \(F/N_\infty\) with the inverse limit of the
actual quotients \(F/N_j\): any compatible system of cosets has
nonempty intersection. The automorphisms just proved commute with
the transition maps, as do their inverses. They therefore induce
an automorphism on \(F/N_\infty\), and
\(\psi_d^{-1}(N_\infty)=N_\infty\).
The original norm gives the exact endomorphism identity
\(\alpha\psi_d=\psi_dc\), without assuming that \(\psi_d\) is
an automorphism of \(F\). Descending this identity and using the
proved quotient automorphisms gives all the stated conjugacies
on the \(F\)-quotients.

For (A.107.3), set \(L=\psi_d(F)\), a closed subgroup.
Surjectivity on \(F/N_\infty\) gives \(F=N_\infty L\), and the exact
preimage equality gives
\[
 L\cap N_\infty
   =\psi_d\bigl(\psi_d^{-1}(N_\infty)\bigr)
   =\psi_d(N_\infty).
 \tag{A.107.8}
\]
These two equalities prove both directions of (A.107.3).
No assertion that the remaining restriction is onto is made.

\emph{The cubic relation is proved before defining the modular map.}
The exact original ordered identity from A.87, equations
(A.87.9)--(A.87.12), with its earlier exponent \(m\) equal to the
present \(k\), is
\[
 \begin{gathered}
 v^2=1,\qquad \alpha(v)=v,\\
 v\delta^d=u^{-1}y^{-k}r,\qquad
 H_d=u\bigl((v\delta^d)^3\bigr)^{-1}u^{-1}.
 \end{gathered}
 \tag{A.107.9}
\]
A.106 therefore implies \((v\delta^d)^3=1\) in \(B_1\).
The universal property of the full free profinite product now
defines an actual continuous homomorphism
\[
 \Psi:B\longrightarrow B_1,\qquad
 s\longmapsto v,\quad r\longmapsto v\delta^d,
 \quad\text{hence }\delta\longmapsto\delta^d .
 \tag{A.107.10}
\]
No homomorphism with target \(B\) has been assumed.
Since \(u\in F\) and \(d\) is odd, (A.107.10) is the identity on the
original quotient \(S_3\). On the actual pure generators it gives
\[
 x\longmapsto x^d,\qquad
 y=sxs\longmapsto u^{-1}y^du,
\]
namely (A.107.2). Its restriction kills \(N_1\) in \(B_1\),
by the already established preservation of \(N_1\). Thus it descends
to an endomorphism \(\Psi_1\) of \(B_1\).
Its kernel lies in \(E_1=F/N_1\), where its restriction is injective.
It is onto by first matching the coset in \(S_3\), then correcting
in \(E_1\). This proves the claimed automorphism.

The original extensions have
\[
 \alpha(s)=s\phi,\quad \alpha(\delta)=\delta^{-1},
 \qquad c_B(s)=s,\quad c_B(\delta)=\delta^{-1}.
\]
Their restrictions preserve \(N_1\), so they descend to \(B_1\).
Using \(\alpha(v)=v\), their values on \(s,\delta\) prove
\(\alpha\Psi=\Psi c_B\). Descent and invertibility of \(\Psi_1\)
give (A.107.5).

For \(j>1\), automorphy on \(F/N_j\) alone does not prove the cubic
relation needed for a map on \(B/N_j\); no assertion \(H_d\in N_j\)
for all \(j\) was used or obtained.
\end{proof}

All maps concern the actual closed subgroups and their quotient
topologies. The residual restriction in (A.107.3), the full hexagon,
and the original full-group recognition problem remain separate.

